To strictly evaluate the logical reasoning capabilities of visual-language models on \BenchmarkName, we adhere to the rigorous evaluation protocols established in M$^3$-Med, formalizing the assessment into two distinct tasks: Temporal Answer Grounding in Single Video and Temporal Answer Grounding in Video Corpus.
% 为了严格评估视觉语言模型在 \BenchmarkName 上的逻辑推理能力，我们遵循 M$^3$-Med 中建立的严格评估协议，将评估形式化为两个不同的任务：单视频时间基础和语料库级时间基础。

\textbf{Task 1: Temporal Answer Grounding in Single Video (TAGSV)}

\textbf{Problem Definition.} Given a natural language query $Q$ and a reference medical video $V$, the objective is to localize the specific temporal segment $T^* = [t_{start}^*, t_{end}^*]$ that semantically answers $Q$.
% 问题定义。给定一个自然语言查询 $Q$ 和一个参考医学视频 $V$，目标是定位语义上回答 $Q$ 的特定时间段 $T^* = [t_{start}^*, t_{end}^*]$。

\textbf{Metric Formulation.} Let $T_i$ denote the predicted interval and $T_i^*$ the ground-truth interval for the $i$-th query. We quantify the localization accuracy using the Temporal Intersection over Union (IoU):
% \textbf{度量公式。}令 $T_i$ 表示第 $i$ 个查询的预测区间，$T_i^*$ 表示第 $i$ 个查询的真实区间。我们使用时间交并比 (IoU) 来量化定位精度：
\begin{align}
    \text{IoU}(T_i, T_i^*) = \frac{\text{duration}(T_i \cap T_i^*)}{\text{duration}(T_i \cup T_i^*)}
\end{align}
Performance is aggregated using 
% \textbf{Mean IoU (mIoU)} and 
\textbf{Recall@$\tau$} (the percentage of queries where IoU exceeds a threshold $\tau \in \{0.3, 0.5, 0.7\}$):
% 性能采用平均 IoU (mIoU) 和召回率@$\tau$（IoU 超过阈值 $\tau \in \{0.3, 0.5, 0.7\}$ 的查询百分比）进行汇总：
\begin{align}
    \text{R@}\tau &= \frac{1}{N} \sum_{i=1}^N \mathbb{I}[\text{tIoU}(T_i, T_i^*) \ge \tau] 
    % \\
    % \text{mIoU} &= \frac{1}{N} \sum_{i=1}^N \text{tIoU}(T_i, T_i^*)
\end{align}
where $N$ is the total number of samples and $\mathbb{I}[\cdot]$ is the indicator function.
% 其中 $N$ 为样本总数，$\mathbb{I}[\cdot]$ 为指示函数。

\textbf{Task 2: Temporal Answer Grounding in Video Corpus (TAGVC)}

\textbf{Problem Definition.} This task extends TAGSV to a retrieval setting. Given a query $Q$ and a video corpus $\mathcal{V} = \{V_1, \dots, V_M\}$, the model must identify the target video $V^*$ and localize the evidence segment $T^*$. The model outputs a ranked list of $K=50$ hypotheses $\{(V_{i,j}, T_{i,j})\}_{j=1}^K$, ordered by confidence.
% 问题定义：本任务将 TAGSV 扩展到检索场景。给定查询 $Q$ 和视频语料库 $\mathcal{V} = \{V_1, \dots, V_M\}$，模型必须识别目标视频 $V^*$ 并定位证据片段 $T^*$。模型输出一个按置信度排序的 $K=50$ 个假设列表 $\{(V_{i,j}, T_{i,j})\}_{j=1}^K$。

\textbf{Metric Formulation.} We define a scoring function $S_{i,j}$ that penalizes incorrect video retrieval:
% \textbf{指标公式} 我们定义一个评分函数 $S_{i,j}$，用于惩罚错误的视频检索：
\begin{align}
    S_{i,j} = \mathbb{I}[V_{i,j} = V_i^*] \cdot \text{tIoU}(T_{i,j}, T_i^*)
\end{align}
To evaluate the combined retrieval and localization performance, we employ the \textbf{Top-$K$ mIoU} metric, which considers the best-localized hypothesis within the top-$K$ retrieved candidates:
% 为了评估检索和定位的综合性能，我们采用 \textbf{Top-$K$ mIoU} 指标，该指标考虑检索到的前 $K$ 个候选结果中定位最佳的假设：
\begin{align}
    \text{Top-}K\text{ mIoU} &= \frac{1}{N} \sum_{i=1}^{N} \max_{j \le K} S_{i,j}
\end{align}
We report results for $K \in \{1, 10, 50\}$.
% and their aggregate average.
% 我们报告 $K \in \{1, 10, 50\}$ 的结果及其平均值。

\textbf{Baselines and Comparative Analysis}
To benchmark the complexity of \BenchmarkName, we evaluate three distinct tiers of solvers:
% 为了衡量 \BenchmarkName 的复杂度，我们评估了三个不同级别的求解器：
\begin{itemize}
    \item \textbf{Supervised Specialist:} MutualSL \cite{mutualsl}, a baseline trained specifically on medical video grounding tasks.
    \item \textbf{Heuristic Baseline} (Silver Standard): A fine-tuned CLIP-based matcher, representing a purely semantic approach without structural reasoning capabilities.
    \item \textbf{Chain-of-Thought Oracle} (Golden Standard): GLM-4.5V \cite{glm}, a state-of-the-art MLLM equipped with high-resolution vision and long-context window. We utilize a specific "Step-by-Step" prompting strategy to elicit explicit reasoning chains, serving as an upper-bound proxy for reasoning capability.
\end{itemize}
% \item \textbf{监督专家：} MutualSL \cite{mutualsl}，一个专门针对医学视频定位任务训练的基线模型。
% \item \textbf{启发式基线}（银级标准）：一个经过微调的基于 CLIP 的匹配器，代表一种纯语义方法，不具备结构推理能力。
% \item \textbf{思维链预言机}（金级标准）：GLM-4.5V \cite{glm}，一个配备高分辨率视觉和长上下文窗口的先进多层逻辑模型。我们采用一种特定的“逐步”提示策略来引出明确的推理链，作为推理能力的上限代理。

Beyond standard metrics, we conduct a \textbf{Logic Alignment Analysis} for the Oracle baseline. We quantify the correlation between the \textit{Prescribed Hop Count} (defined by our generation pipeline) and the \textit{Inferred Hop Count} (the number of reasoning steps actually emitted by GLM-4.5V).
% 除了标准指标之外，我们还对 Oracle 基线进行了逻辑对齐分析。我们量化了预设跳数（由我们的生成管道定义）与推断跳数（GLM-4.5V 实际发出的推理步骤数）之间的相关性。

\textbf{Experimental Results and Observations} 
The experimental outcomes validate following critical hypotheses regarding \BenchmarkName:
% 分析。实验结果验证了关于基准测试的两个关键假设：
\begin{itemize}
    \item \textbf{Human-Parity Complexity}: The performance gap between the Supervised Specialist and the Oracle indicates that our synthesized queries achieve a level of complexity comparable to expert-curated problems, necessitating advanced reasoning beyond simple pattern matching.
    \item \textbf{Distinguishability of Logic}: The performance stratification between Simple and Complex queries is pronounced across all models. Notably, the CLIP-based Heuristic fails significantly on Complex queries (Multi-hop), whereas the CoT Oracle maintains robustness.
    \item \textbf{Neuro-Symbolic Alignment}: As illustrated in the alignment heatmap (Fig. \ref{fig:heatmap}), there exists a strong positive correlation between the pipeline's prescribed complexity and the Oracle's reasoning depth. This confirms that our automated pipeline successfully encodes valid, deterministically solvable multi-hop logic into the generated benchmarks.
\end{itemize}
% \item \textbf{人类水平复杂度}：监督专家模型与预言机之间的性能差距表明，我们合成的查询达到了与专家精心挑选的问题相当的复杂度水平，这需要超越简单模式匹配的高级推理能力。
% \item \textbf{逻辑可区分性}：所有模型在简单查询和复杂查询之间的性能差异都非常显著。值得注意的是，基于 CLIP 的启发式算法在复杂查询（多跳）上表现明显不佳，而 CoT 预言机则保持了良好的鲁棒性。
% \item \textbf{神经符号对齐}：如图 X 所示的对齐热图所示，流水线预设的复杂度与预言机的推理深度之间存在很强的正相关性。这证实了我们的自动化流水线能够成功地将有效的、确定性可解的多跳逻辑编码到生成的基准测试中。
